\section{NTK recursion}

We now state the main result characterizing the NTK recursion for MMNNs. This recursion governs how the kernel evolves through layers and determines the training dynamics in the infinite-width limit.

\begin{theorem}[Infinite-width NTK recursion]\label{thm:ntk_recursion}
Under NTK parameterization \( 1/\sqrt{\bm{n}_\ell} \) and a sequential limit \( \bm{n}_1\to\infty,\dots,\bm{n}_L\to\infty \) with \( \bm{w}^{(\ell)} \sim \mathcal{N}(0, \mathbf{I}_{\bm{d}_{\ell-1}}) \) i.i.d. frozen after initialization, the layer-\( \ell \) NTK satisfies
\begin{align}
\Kop^{(0)}(x,x') \;&=\; 0, \\
\Kop^{(1)}(x,x') \;&=\; 1 + \Sigma^{(1)}(x,x'), \\
\Kop^{(\ell)}(x,x') \;&=\; 1 + \frac{1}{r}\,\Kop^{(\ell-1)}(x,x')\cdot\dot{\Sigma}^{(\ell)}(x,x')\;+\;\frac{1}{r}\,\Sigma^{(\ell)}(x,x'),\quad \ell\ge2,
\end{align}
where the factor \( 1/r \) appears for bottleneck layers \( \ell \ge 2 \) (low-rank regime). For \( \ell=1 \), \( \Sigma^{(1)} \) is deterministic; for \( \ell>1 \), \( \Sigma^{(\ell)} \) and \( \dot{\Sigma}^{(\ell)} \) are random fields whose fluctuations shrink with the rank \( r \) of bottlenecks.

\vspace{0.2cm}
Furthermore, the gradient flow dynamics in the NTK regime satisfy
\[
\frac{d}{dt}f_t(x) = -\int \Kop^{(L)}(x,x') \big(f_t(x') - y(x')\big) d\mu(x'),
\]
where \( \mu \) is the data distribution and \( y \) is the target function.
\end{theorem}

Proof: See Appendix~\ref{appendix:proof-ntk-recursion}, Proof~\ref{proof:ntk_recursion}.

\vspace{0.3cm}
\paragraph{No weights particle coupling}
We can now interpret the structure of the recursion. 
Similarly to the original MLPs NTK recursion \cite{jacot2018neural}, the additive term \( \Sigma^{(\ell)} \) encodes fresh basis contribution at layer \( \ell \), while the multiplicative coupling \( \Kop^{(\ell-1)}\cdot\dot{\Sigma}^{(\ell)} \) reflects the backward through layers via the derivative kernel, since there is no coupling between weights particles between 2 layers during training due to random features in the middle. 
Nonetheless, this decoupling that makes the training easier to tackled under the mean-field limit is effectively lost at the NTK limit and allows to factorize the derivative kernel.

\vspace{0.4cm}
We can now derive an explicit closed-form expression for the NTK at any depth by expanding the recursion.

\begin{corollary}[General L-layer NTK; closed form]\label{cor:general_L_layer_ntk}
Under the assumptions of Theorem~\ref{thm:ntk_recursion}, for any \( L\ge1 \), the NTK admits the explicit form:
\begin{equation}
\label{eq:ntk_explicit_form}
\Kop^{(L)}(x,x') = \sum_{\ell=1}^{L} \frac{1}{r^{L-\ell}}\left(\prod_{k=\ell+1}^{L} \dot{\Sigma}^{(k)}(x,x')\right) + \sum_{\ell=1}^{L} \frac{1}{r^{\max(1,L-\ell)}}\,\Sigma^{(\ell)}(x,x') \prod_{k=\ell+1}^{L} \dot{\Sigma}^{(k)}(x,x'),
\end{equation}
where the first sum has bias contributions (empty product \( = 1 \) when \( \ell=L \)) and the second has fresh-basis contributions; each bottleneck layer contributes a factor \( 1/r \). For \( L=2 \): \( \Kop^{(2)} = 1 + \frac{1}{r}(\dot{\Sigma}^{(2)} + \Sigma^{(1)}\dot{\Sigma}^{(2)} + \Sigma^{(2)}) \).
\end{corollary}

\begin{remark}
Equivalently, expanding the recursion yields a quadratic number of terms. For example, for \( L=3 \):
\begin{align}
\Kop^{(3)}(x,x') &= 1 + \frac{1}{r}\dot{\Sigma}^{(2)}\cdot\dot{\Sigma}^{(3)} + \frac{1}{r}\dot{\Sigma}^{(3)} \nonumber\\
&\quad + \frac{1}{r^2}\Sigma^{(1)}\cdot\dot{\Sigma}^{(2)}\cdot\dot{\Sigma}^{(3)} + \frac{1}{r}\Sigma^{(2)}\cdot\dot{\Sigma}^{(3)} + \frac{1}{r}\Sigma^{(3)}.
\end{align}
Each bottleneck layer contributes a factor \( 1/r \); layer 1 has no such factor. In total, there are \( 2L \) terms for depth \( L \).
\end{remark}

\begin{proof}
\label{proof:general_L_ntk_inline}
By expanding the recursion in Theorem~\ref{thm:ntk_recursion}. See Appendix~\ref{appendix:proof-general-L} for the detailed proof by induction.
\end{proof}

\vspace{0.3cm}
\paragraph{Structure at depth \( L \).}
Iterating the recursion yields a quadratic number of terms. From the explicit form \eqref{eq:ntk_explicit_form}, we see two types of contributions:
\begin{itemize}
    \item \textbf{Bias terms:} Each layer's \( +1 \) contribution multiplies by \( 1/r^{L-\ell} \) (for bottleneck layers) and through all subsequent derivative kernels, giving \( L \) terms.
    \item \textbf{Fresh basis terms:} Each layer's \( \Sigma^{(\ell)} \) contribution carries \( 1/r^{L-\ell} \) and propagates through subsequent derivative kernels, giving \( L \) terms.
\end{itemize}
In total, the expansion contains \( 2L \) terms, each representing a different path through the network architecture. This linear growth in the number of terms (with respect to depth) becomes quadratic when considering all combinations of terms in finite-width corrections, making explicit computation and analysis very difficult to handle for deep networks.

\vspace{0.4cm}
\subsection{Large-depth regime: probabilistic kernel and exponential decay}
\label{subsec:large-depth}

We now outline how to adapt the recursion for large depth \( L \), focusing on the probabilistic (mean) kernel and the role of the \( 1/r \) scaling at low-rank bottlenecks.

\paragraph{Probabilistic kernel recursion and GP concatenation.}
For \( \ell > 1 \), \( \Sigma^{(\ell)} \) and \( \dot{\Sigma}^{(\ell)} \) are random fields because they depend on \( \bm{h}^{(\ell-1)} \), which is the output of a (conditional) Gaussian process. By Theorem~\ref{thm:gp-composition}, the hidden states form a \emph{composition of GPs}: each layer's output is conditionally Gaussian given the previous layer, but the unconditional law is a deep GP. The concatenation structure implies that at layer \( \ell \), the pair \( (\bm{h}^{(\ell-1)}(x), \bm{h}^{(\ell-1)}(x')) \) follows a \( 2r \)-dimensional Gaussian (in the bottleneck regime), with covariance determined by \( \Sigma^{(\ell-1)} \). The base kernel \( \Sigma^{(\ell)} \) and derivative kernel \( \dot{\Sigma}^{(\ell)} \) are thus expectations over \( \bm{w}^{(\ell)} \), conditioned on this Gaussian pair; they inherit randomness from \( \Sigma^{(\ell-1)} \).

The \emph{mean} (probabilistic) NTK \( \tilde{\Kop}^{(L)} = \mathbb{E}[\Kop^{(L)}] \) satisfies a recursion involving expectations of products:
\begin{equation}
\label{eq:mean_ntk_recursion}
\tilde{\Kop}^{(\ell)} \;=\; 1 \;+\; \mathbb{E}\!\big[\Kop^{(\ell-1)}\cdot\dot{\Sigma}^{(\ell)}\big] \;+\; \mathbb{E}\!\big[\Sigma^{(\ell)}\big],\quad \ell \ge 1.
\end{equation}
Since \( \Kop^{(\ell-1)} \) and \( \Sigma^{(\ell)}, \dot{\Sigma}^{(\ell)} \) are not independent (they share the randomness of \( \bm{h}^{(\ell-1)} \)), we have \( \mathbb{E}[\Kop^{(\ell-1)}\cdot\dot{\Sigma}^{(\ell)}] \neq \tilde{\Kop}^{(\ell-1)}\cdot\mathbb{E}[\dot{\Sigma}^{(\ell)}] \) in general. The correct expansion requires tracking covariances between terms across layers. Under Fisher--Kibble decoupling at each layer, the randomness decomposes into independent angular and radial parts, which can be handled iteratively; the technical difficulty lies in controlling the nested expectations.

\paragraph{\( 1/r \) scaling at each bottleneck.}
In the low-rank regime, the derivative kernel at layer \( \ell \ge 2 \) scales with \( 1/r \). For ReLU and isotropic weights with \( \mathrm{Cov}(w) = I_r/r \) on the \( r \)-dimensional bottleneck output, one has
\[
\dot{\Sigma}^{(\ell)}(x,x') \;=\; \frac{1}{r}\,\mathbb{E}_w\!\Big[\dot{\sigma}\!\big(w^\top \bm{h}^{(\ell-1)}(x)\big)\,\dot{\sigma}\!\big(w^\top \bm{h}^{(\ell-1)}(x')\big)\,\|w\|^2\Big] \;=\; \frac{1}{r}\,\bar{\dot{\Sigma}}^{(\ell)}(\rho_\ell),
\]
where \( \rho_\ell = \cos\angle(\bm{h}^{(\ell-1)}(x), \bm{h}^{(\ell-1)}(x')) \) and
\(
\bar{\dot{\Sigma}}^{(\ell)}(\rho)
\)
denotes the corresponding scalar derivative kernel evaluated at cosine similarity \(\rho\in[-1,1]\).
For ReLU under the EOC normalization \(\mathrm{Tr}(\mathrm{Cov}(w))=1\), the derivative kernel has the explicit form
\[
\bar{\dot{\Sigma}}^{(\ell)}(\rho)=\mathbb{P}(Z>0,Z'>0)=\frac{1}{2}-\frac{\arccos(\rho)}{2\pi}\in\Big[0,\frac12\Big],
\]
where \((Z,Z')\) is centered Gaussian with correlation \(\rho\). In particular, \(\sup_{\rho\in[-1,1]}\bar{\dot{\Sigma}}^{(\ell)}(\rho)=1/2\), uniformly in \(\ell\).
Similarly, \( \Sigma^{(\ell)} \) involves a norm product \( \|\bm{h}^{(\ell-1)}\|^2/r \) with \( \mathbb{E}[\cdot] \to 1 \) and \( \mathrm{Var}(\cdot) = O(1/r) \). Thus, at each bottleneck layer, both the multiplicative and additive contributions carry a \( 1/r \) factor relative to the full-width case.

\paragraph{Exponential decay with depth.}
Assume for simplicity that all layers \( \ell \ge 2 \) are bottlenecks with the same rank \( r \). For ReLU under EOC we may take the explicit uniform bound
\[
\bar{\dot{\Sigma}}^{(\ell)}(\rho)\le c_0:=\frac12,\qquad \forall \rho\in[-1,1],\ \forall \ell\ge 2.
\]
From the closed form \eqref{eq:ntk_explicit_form}, the contribution from layer \( \ell \) to \( \Kop^{(L)} \) is multiplied by \( \prod_{k=\ell+1}^{L} \dot{\Sigma}^{(k)} \). With \( \dot{\Sigma}^{(k)} = \bar{\dot{\Sigma}}^{(k)}(\rho_k)/r \), we obtain
\[
\prod_{k=\ell+1}^{L} \dot{\Sigma}^{(k)} \;\sim\; \frac{1}{r^{L-\ell}}\,\prod_{k=\ell+1}^{L} \bar{\dot{\Sigma}}^{(k)}(\rho_k) \;\le\; \frac{c_0^{L-\ell}}{r^{L-\ell}}.
\]
Hence, contributions from early layers (\( \ell \ll L \)) decay \emph{exponentially} in \( L - \ell \) when \( c_0/r < 1 \). The dominant terms in \( \Kop^{(L)} \) come from the last \( O(1) \) layers; the NTK becomes increasingly localized to the top of the network as depth grows.

\begin{proposition}[Informal: depth-induced decay]\label{prop:depth_decay}
Under the assumptions above, for \( r \) fixed and \( L \to \infty \), early-layer contributions to \( \Kop^{(L)} \) are suppressed by a factor \( (c/r)^{L-\ell} \) for some \( c = O(1) \). In particular, when \( c/r < 1 \), the NTK is dominated by contributions from the last \( O(\log r) \) layers. The mean NTK \( \tilde{\Kop}^{(L)} \) concentrates around a limit that depends primarily on these top layers.
\end{proposition}

\paragraph{Open directions.}
A full treatment requires: (i)~a rigorous recursion for \( \tilde{\Kop}^{(\ell)} \) with explicit covariance bounds; (ii)~verification that \( \bar{\dot{\Sigma}}^{(\ell)} \) remains uniformly bounded in \( \ell \) and \( \rho_\ell \); (iii)~extension of Fisher--Kibble decoupling to nested layers to control \( \mathbb{E}[\Kop^{(\ell-1)}\dot{\Sigma}^{(\ell)}] \). The exponential decay suggests that very deep MMNNs may exhibit a ``effective depth'' phenomenon similar to ResNets, where only the top layers meaningfully contribute to the NTK.

\vspace{0.4cm}
We now specialize to the three-layer case, where we can obtain a particularly compact form under a homogeneity assumption on the activation function.

\begin{corollary}[three-layer NTK and homogeneous activation compact form]\label{cor:two_layer_ntk}
Under the assumptions of Theorem~\ref{thm:ntk_recursion} with \( L=2 \),
\[
\Kop^{(1)}(x,x') \;=\; 1 + \Sigma^{(1)}(x,x'),
\]
\[
\Kop^{(2)}(x,x') \;=\; 1 + \frac{1}{r}\,\Kop^{(1)}(x,x')\cdot\dot{\Sigma}^{(2)}(x,x') \;+\; \frac{1}{r}\,\Sigma^{(2)}(x,x').
\]
\end{corollary}

\begin{proof}
\label{proof:two_layer_ntk_inline}
Immediate from Theorem~\ref{thm:ntk_recursion}.
\end{proof}




\vspace{0.3cm}
We can now state a more compact representation under a homogeneity assumption. 
\vspace{0.3cm}


\subsection{Microscopic distributions: Fisher–Kibble decoupling}
\begin{lemma}[Fisher–Kibble decoupling]\label{lem:fisher-kibble}
Let \(x,y\) be input vectors and let \(x_1,y_1\) denote their rank-\(r\) random projections. Define the empirical correlation \( \rho_1 = \cos\angle(x_1,y_1) \) and squared norms \( u=\|x_1\|^2, v=\|y_1\|^2 \). Then:
\begin{itemize}
    \item \( \rho_1 \) follows Fisher's correlation distribution~\cite{fisher1915probable,hotelling1953new}, centered at the true correlation \( \rho \).
    \item \( (u,v) \) follow Kibble's (1941) bivariate Gamma (chi-square) law.
    \item Angular and radial parts are independent: \( p(\rho_1,u,v) = p_{\text{Fisher}}(\rho_1)\,p_{\text{Kibble}}(u,v) \).
\end{itemize}
\end{lemma}

Proof: See Appendix~\ref{appendix:fisher-kibble}, Proof~\ref{proof:fisher_kibble}.

\paragraph{Homogeneity assumption enabling decoupling.}
The independence of angular and radial components relies on the \emph{positive 1-homogeneity} of the ReLU activation: \( \ReLU(\alpha z) = \alpha \ReLU(z) \) for \( \alpha > 0 \). This property ensures that the base kernel factorizes as
\[
\Sigma^{(1)}(x,x') = \|x\|\,\|x'\| \cdot \tilde{\Sigma}^{(1)}(\rho),
\]
where \( \tilde{\Sigma}^{(1)}(\rho) \) depends only on the cosine similarity \( \rho = \langle x,x'\rangle/(\|x\|\|x'\|) \), enabling the separation of radial norms from angular correlations. For isotropic Gaussian projections, rotational invariance guarantees that the empirical correlation \( \rho_1 \) and the norms \( (\|x_1\|,\|y_1\|) \) are statistically independent, yielding the factorization stated in the lemma.

This decoupling result enables us to derive the asymptotic distribution of the empirical kernel. We can state the following central limit theorem.



In the three-layer low-rank setting, denote
\[
\rho_1 \;=\; \cos\angle(\bm{h}^{(1)}(x),\bm{h}^{(1)}(x')),\qquad w_r \;=\; \frac{\|\bm{h}^{(1)}(x)\|\,\|\bm{h}^{(1)}(x')\|}{r},
\]
then \emph{both \(\rho_1\) and \(w_r\) are random variables} induced by the randomness of the features and the rank-
\(r\) projection; their values fluctuate with the initialization and with \(r\).
The NTK admits the compact representation
\[
\Kop^{(2)}(x,y) \;=\; 1 \;+\; \frac{1}{r}\,\Kop^{(1)}(\rho_1)\!\left(1-\frac{\arccos(\rho_1)}{\pi}\right) \;+\; \frac{1}{r}\,w_r\,\Sigma^{(1)}(\rho_1)
\]
with \( \Sigma^{(1)}(u) = \frac{1}{\pi}\big(\sqrt{1-u^2}+u(1-\arccos u)\big) \).

\vspace{0.3cm}
\paragraph{Remark (Fisher–Kibble decoupling and independence).}
\label{rmk:fisher-kibble-law}
We can now characterize the statistical properties of the empirical correlation and norms. By rotational invariance, the empirical correlation and squared norms follow Fisher–Kibble decoupling~\cite{fisher1915probable,hotelling1953new}:
\[
\rho_1 \sim \text{Fisher}(\rho,r),\qquad (\|\bm{h}^{(1)}(x)\|^2,\|\bm{h}^{(1)}(x')\|^2)\sim \text{Kibble}(r,\rho),
\]
 

\vspace{0.2cm}
We now state the explicit forms of these distributions. The \textbf{Fisher distribution density}~\cite{fisher1915probable} is
\[
\boxed{p(\rho_1 \mid \rho, r) = \frac{(r-2)\,\Gamma(r-1)\,(1-\rho^2)^{\frac{r-1}{2}}\,(1-\rho_1^2)^{\frac{r-4}{2}}}{\sqrt{2\pi}\,\Gamma(r-\tfrac{1}{2})\,(1-\rho\rho_1)^{r-\frac{3}{2}}}\,{}_2F_1\!\left(\tfrac{1}{2},\tfrac{1}{2};\, r-\tfrac{1}{2};\, \tfrac{1+\rho\rho_1}{2}\right),}
\]
for \( r > 2 \). For large \( r \), the inverse Gudermann function applied to the complementary angle follows a normal distribution:
\[
\text{gd}^{-1}\!\left(\frac{\pi}{2} - \arccos(\rho_1)\right) \;\approx\; \mathcal{N}\!\left(\text{arctanh}(\rho),\; \frac{1}{r-3}\right)
\]
where \( \rho_1 = \cos(\theta_1) \) is the empirical correlation and \( \theta_1 \) is the angle between the projected vectors. In particular, this implies \( \text{Var}(\rho_1) = O(1/r) \).

\vspace{0.2cm}
Similarly, \textbf{Kibble's~\cite{kibble1941two} bivariate chi-square law} for \( (u,v) = (\|\bm{h}^{(1)}(x)\|^2, \|\bm{h}^{(1)}(x')\|^2) \) has density
\[
\boxed{f(u,v) = \frac{(uv)^{\frac{r/2-1}{2}}\,\exp\!\left(-\tfrac{u+v}{2(1-\rho^2)}\right)}{\Gamma(r/2)\,\big(2(1-\rho^2)\big)^{\frac{r}{2}+1}\,\rho^{\frac{r/2-1}{2}}}\, I_{\frac{r}{2}-1}\!\left(\frac{\rho\sqrt{uv}}{1-\rho^2}\right),\qquad u,v \ge 0,}
\]
where \( I_\nu \) is the modified Bessel function of the first kind. Defining \( w_r = \sqrt{uv}/r \), concentration results (see Theorem~\ref{thm:rkhs_rank_concentration}) imply \( \mathbb{E}[w_r] \to 1 \) and \( \mathrm{Var}(w_r)=O(1/r) \).

\vspace{0.15cm}
The moments of this distribution are given by:
\[
\mathbb{E}[u]=\mathbb{E}[v]=r,\quad \text{Var}(u)=\text{Var}(v)=2r,\quad \text{Cov}(u,v)=2r\,\rho^2.
\]
 
The derivation of these distributions from Gaussian projections and the independence property are given in Appendix~\ref{appendix:fisher-kibble}.
